\section{Autonomous Landing Proof of Concept}
\label{section:proof_of_concept}

The last part of the initial phase of this research has been to successfully develop a proof of concept
of the autonomous landing algorithm on a physical drone, which was also the goal of the first part (in Section \ref{section:initial_hexacopters}).
The crucial difference is that we have used a DJI Spark as the drone platform,
which is stable enough to remain nearly still without human intervention,
even in a GPS-denied environment.
However, even though the proof of concept is successful,
its implementation is significantly constrained because of its basis on the DJI Mobile SDK,
as explained below.
The DJI Mobile SDK is necessary because the DJI Spark has no companion board,
but anticipated future drone systems will have a more ideal setup using the DJI Onboard SDK.

\subsection{System Architecture}

The system is outlined in Figure~\ref{figure:spark_architecture}.
All data is transmitted from the Spark to its controller via a wireless connection.
The controller interfaces with an Android tablet to expose the drone's data and API.
The app must authenticate with DJI using a user-specific API key
before it can communicate with the drone,
which drives home the proprietary, block-box nature of the system.
Importantly, after the app has successfully authenticated once, it can re-authenticate without
an internet connection for some period of time,
meaning that an internet connection is not required in the field.
The Android app then receives data about the condition of the drone,
including all flight data (e.g.~attitude, velocities, flight modes, etc.),
and the video stream from the camera.
The app periodically decodes bitmap frames from the video stream,
compresses them into black-and-white .webp images,
and transmits them over a WiFi connection to a companion board.
The companion board then carries detects April Tags in the video frames,
determines what action to take according to the control policy,
instantiates a command message,
and sends that command back to the app.
The app then reads the command,
creates a virtual stick input via the DJI SDK,
updates the GUI to visually show the command,
and sends the command to the drone.
The drone then carries out the command.

\begin{figure}[ht]
    \centering
    \includegraphics[width=0.9\textwidth]{images/spark_architecture.drawio}
    \caption{Data flow for the Spark autonomous landing system.}
    \label{figure:spark_architecture}
\end{figure}

The virtual stick inputs simulate input as if it were coming from the controller.
Each stick command can be in the interval $[-1, 1]$.
The drone can interpret those stick commands as setpoints either for position or velocity.
Since we are abstracting from position, we only use the velocity mode.
Therefore, a command of $1$ on the pitch stick
means that the drone should move forward at its maximum velocity.

\subsection{Results}

The system is able to reliably land the drone without human intervention.
\footnote{Autonomous landing demonstration video: \url{https://vimeo.com/manage/videos/646612738}}
So far, we have tested it with April Tag 48h12 and April Tag 24h10 landing pads.
As anticipated, the April Tag 24h10 tags are subject to somewhat common instances of orientation ambiguity,
whereas the April Tag 48h12 tags are mostly free of orientation ambiguity.
In the demonstration video, this orientation ambiguity appears as jittery control inputs from the companion board,
as well as (sometimes) jittery behavior from the drone.
Surprisingly - and fortunately - this behavior appears to be non-destructive in this case.

The Android app has 5 sliders to control the drone:
\begin{enumerate}
    \item Pitch (translated to forward/backward speed)
    \item Roll (translated to left/right speed)
    \item Yaw (translated to yaw speed, i.e.~rotational speed in the Z axis)
    \item Throttle (translated to vertical speed)
    \item Gimbal Tilt (translated to gimbal tilt speed)
\end{enumerate}
As well as buttons to enable or disable virtual control in the flight controller,
to command the drone to takeoff,
and to command the drone to land (blindly).
The app works just well enough to allow for basic testing of this proof of concept,
but has only received minimal development efforts because it is almost pure overhead.
Therefore, it crashes often.
The proposed research will avoid the Mobile SDK and will therefore avoid the need for a custom app.
However, subsequent app development can be the subject of a future project if necessary.
Still, the app has worked well enough to allow for many autonomous landings.
\footnote{The autonomous landing demonstration from the point of view of the app: \url{https://vimeo.com/manage/videos/646620944}}

\subsection{Limitations and Implementation Quirks}

The Mobile SDK architecture is not well-suited to this application for a number of reasons.
First, there is only one wired connection on the tablet, meaning that at least one connection
must be wireless - either the connection to the controller, or the connection to the companion board.
The tablet is not fast enough to do April Tag image analysis on its own,
and there would be significantly increased development efforts required to install April Tag
on the tablet alone (so the companion board is required).
Moreover, connecting the tablet to the controller via its WiFi interface incurs prohibitive latency
as a result of the large amount of processing required.
It is also prohibitively slow to connect two clients over the controller's WiFi simultaneously.
Second, the frames of the video stream are available in bitmap format,
and each frame is prohibitively large/slow to transmit over a wireless connection to the companion board.
In order to transmit ``video'' in ``real time,'' the app converts bitmap frames to black and white,
and then compresses them into .webp format with 10\% quality.
This allows a framerate of about 6 Hz, with a delay in each frame of about 0.25 seconds,
to be delivered to the companion board.
Amazingly, this is enough to land the drone with reasonable accuracy.
Third, the Spark requires a blind commitment in the last stage of landing
- where the flight controller forces the gimbal to point directly forward,
apparently to protect it from the ground since it has very close clearance.
So, although the drone responds to commands during this phase of landing,
the companion board cannot see the landing pad and therefore cannot generate commands
for last minute correction.